\fancyhead{}
\fancyfoot{}
\lhead{Introducción}
\cfoot{\thepage}

\chapter{Introducción}

El registro de asistencia es una tarea clave en la gestión educativa y administrativa de instituciones académicas como la 
Facultad Politécnica de la Universidad Nacional del Este (FPUNE). Sin embargo, los métodos tradicionales de control 
manual mediante listas suelen generar errores humanos que afectan la precisión de los registros y la toma de decisiones 
académicas. Este Trabajo Final de Grado (TFG) tiene como objetivo desarrollar un sistema automatizado basado en reconocimiento facial para optimizar el proceso de registro de asistencia en la FPUNE. 
Para ello, se emplearán algoritmos avanzados de biometría que analicen imágenes en tiempo real, asegurando una identificación precisa y eficiente. La metodología contempla la selección tecnológica, el procesamiento de datos, el entrenamiento de modelos y pruebas piloto en un entorno académico durante el año 2025, buscando reducir errores y mejorar la gestión administrativa. De esta forma, el estudio pretende aportar una herramienta que brinde datos más confiables y contribuya a una mayor eficiencia operativa, beneficiando a estudiantes, docentes y autoridades mediante la adopción de tecnologías innovadoras que fortalezcan la gestión académica.

\begin{table}[]
\begin{tabular}{lllll}
\end{tabular}
\end{table}
\section{Motivación}

En la actualidad, el control de asistencia en instituciones educativas sigue dependiendo de métodos manuales, como el llamado de lista o el registro en papel, lo que puede generar errores humanos y pérdida de tiempo en la gestión académica. La automatización de este procedimiento mediante tecnologías de reconocimiento facial se presenta como una alternativa innovadora que no solo mejora la precisión del registro de asistencia, sino que también permite una gestión más eficiente de la información en tiempo real.

En el ámbito académico, este estudio representa una oportunidad para explorar el uso del reconocimiento facial en entornos educativos, impulsando la incorporación de tecnologías modernas que mejoren la gestión institucional. Al haber demostrado eficacia en distintos sectores, su adaptación al contexto universitario puede marcar una diferencia significativa en la optimización de procesos y en la precisión del control de asistencia.

Este proyecto también busca reducir el impacto ambiental al eliminar el uso de papel en los registros de asistencia, promoviendo prácticas más sostenibles dentro de la institución. En este sentido, la motivación principal radica en la combinación de innovación tecnológica, eficiencia administrativa y sostenibilidad, factores clave que justifican la importancia de esta investigación.

\section{Definición del problema}


En la FPUNE, el control de asistencia aún utiliza métodos tradicionales, como el llamado de lista o el registro manual en papel. Estos métodos generan frecuentemente errores humanos, pérdida o deterioro de documentos y dificultades en la organización y almacenamiento eficiente de los datos recopilados. 

Además, estas fallas afectan negativamente la capacidad de las autoridades académicas y administrativas para detectar patrones de asistencia, identificar ausencias repetitivas y aplicar medidas correctivas oportunas.

Ante esta problemática, surge la necesidad de desarrollar una solución tecnológica basada en reconocimiento facial que permita registrar automáticamente la asistencia.
\section{Objetivos}
\subsection{Objetivo General}
Desarrollar un prototipo de sistema de reconocimiento facial en un aula para el control de asistencia de estudiantes de la Facultad Politécnica de la Universidad Nacional del Este.
\subsection{Objetivos Específicos}
\begin{enumerate}
    \item Seleccionar los algoritmos de reconocimiento facial.
    \item Seleccionar las herramientas tecnológicas para el desarrollo del Sistema.
    \item Identificar los requisitos para el desarrollo del sistema de registro de asistencia para la FPUNE.
    \item Desarrollar el sistema de reconocimiento facial y gestión de asistencias.
    \item Realizar pruebas del funcionamiento del prototipo
    \item Evaluar resultados obtenidos.
\end{enumerate}


\section*{Preguntas de Investigación}
\subsection*{Pregunta General.}
¿Cómo desarrollar un sistema de Reconocimiento Facial para el control de asistencia de los estudiantes de la FPUNE?

\subsection*{Preguntas Específicas.}
\begin{enumerate}
    \item ¿Cuáles son los algoritmos de reconocimiento facial?
    \item ¿Cuáles son las herramientas tecnológicas para el desarrollo del Sistema?
    \item ¿Cuáles son los requerimientos para el registro de asistencia de la Facultad Politécnica?
    \item ¿Cómo se desarrollará el sistema de reconocimiento facial y gestión de asistencias?
    \item ¿Qué pruebas se deben realizar para validar su correcto funcionamiento?
    \item ¿Cómo se evaluarán los resultados obtenidos para medir la efectividad del sistema?
    \item ¿Qué indicadores o métricas se pueden utilizar para evaluar el desempeño del sistema?
\end{enumerate}


\section{Hipótesis}
Establecer que el sistema de reconocimiento facial basado en los modelos \textit{YOLOv8n-Face} y \textit{ArcFace} que alcanzará una precisión superior al 85\,\% en el registro automatizado de asistencia bajo condiciones de iluminación controlada (\(\geq 300\,\text{lux}\)) en entornos reales de aula, reduciendo en al menos un 70\,\% el tiempo promedio requerido por los métodos manuales tradicionales utilizados en la Facultad Politécnica de la Universidad Nacional del Este.



\section{Fundamentación.}

El desarrollo de un ``Sistema de Asistencia Basado en Reconocimiento Facial para la FPUNE'' responde a la necesidad de modernizar los procesos administrativos en las instituciones educativas y de aprovechar las ventajas que ofrecen las tecnologías de reconocimiento facial. A continuación, se presentan las diversas fundamentaciones en base a los puntos siguientes:

\begin{itemize}
    \item \textbf{Innovación tecnológica:} El desarrollo de un prototipo de sistema de asistencia basado en reconocimiento facial en una institución
educativa como la FPUNE representa una innovación tecnológica que
puede servir como modelo para otras instituciones y/o otros rubros.
    \item \textbf{Incremento de la precisión en los datos de asistencia:} Los datos
obtenidos del sistema serán más precisos y confiables, lo que permitirá
una mejor evaluación del desempeño académico de los estudiantes y la
identificación de patrones de asistencia.
    \item \textbf{Fomento de la puntualidad y asistencia:} El desarrollo del
    sistema puede incentivar a los estudiantes a ser más puntuales y asistir
    con regularidad a clases.
\end{itemize}


\section{Delimitación del Trabajo.}
\begin{enumerate}
    \item \textbf{Límite del contenido:} Desarrollo de un sistema reconocimiento facial
    inteligente para la asistencia de alumnos de la FPUNE.
    \item \textbf{Límite espacial:} Se limitará las pruebas a ser realizadas en un entorno
    controlado y luego a ser probado en un aula de la FPUNE.
    \item \textbf{Límite Temporal:} Pruebas y desarrollo del sistema durante un semestre
    académico en el año 2025. Incluyendo la recopilación de datos, el entrenamiento del modelo y la evaluación de
    su desempeño.
\end{enumerate}


\section*{Impacto de la Investigación.}

El desarrollo del prototipo ``Sistema de Asistencia Basado en Reconocimiento Facial para la FPUNE'' generará un impacto significativo tanto a nivel institucional como en el entorno social más amplio. A nivel institucional, se mejorará la eficiencia y precisión en el registro de asistencia, reduciendo el margen de error humano. Socialmente, el sistema fortalecerá la confianza en los procesos educativos y administrativos de la FPUNE. Además, la digitalización de los procesos de registro de asistencia, al reducir drásticamente el uso de papel y tinta, contribuirá a disminuir el consumo de recursos.
